% assets/w06-berth-phases.svg — 접안 전략 둘: 옆걸음 대 전진
%
%   bash _tools/tikz2svg.sh assets/src/w06-berth-phases.tex assets/w06-berth-phases.svg
%
% 이 그림이 말하는 것
%   같은 선석, 같은 접근점인데 **마지막에 배가 어느 쪽을 보고 서 있느냐**만 다르다.
%   그 하나가 바뀌면 어느 추진기가 일하는지가 통째로 바뀐다.
%
%   (가) 옆걸음 접안   psi_f = psi_t           부두와 나란히 서서 **옆으로** 붙는다
%                                              → 선체 y 방향 운동, 선수추진기가 주력
%   (나) 전진 접안     psi_f = psi_t - s*90도   부두를 보고 **앞으로** 붙는다
%                                              → 선체 x 방향 운동, 좌우 추진기가 주력
%
% 단계 번호는 모델 안의 phase 신호와 같다
%   0 접근 → 1 선수각 정렬 → 2 접안 → 3 정지
%
% 좌표 — w06-berth-geometry.tex 와 같은 체계를 쓴다
%   선석이 패널 원점. \TN{a}{b} = 부두를 따라 a m, 부두 바깥으로 b m
%   여기서는 6주차 실제 설정 psi_t = 0, s = +1 이므로 t = 북, n = 동이다
%   화면 회전각 = 90 - psi  (북이 위, 동이 오른쪽)

\documentclass[border=6pt]{standalone}
% capstone-style.tex — 이 볼트의 TikZ 그림이 공유하는 스타일
%
% 쓰는 법:  % capstone-style.tex — 이 볼트의 TikZ 그림이 공유하는 스타일
%
% 쓰는 법:  % capstone-style.tex — 이 볼트의 TikZ 그림이 공유하는 스타일
%
% 쓰는 법:  \input{capstone-style}   (assets/src/ 안에서)
% 빌드:     bash _tools/tikz2svg.sh assets/src/w02-node-topic.tex assets/w02-node-topic.svg
%
% 왜 TikZ 인가
%   손으로 SVG 좌표를 쓰면 화살촉·선 굵기·글자 크기가 그림마다 미묘하게 달라진다.
%   그것이 "자료가 어설퍼 보인다"의 정체다. 스타일을 한 번만 정하고 상대좌표로
%   배치하면 좌표를 고쳐도 그림이 무너지지 않는다.
%   수식은 본문과 같은 TeX 글꼴로 나온다.
%
% 한글
%   ko.TeX(DVI 경로)를 쓴다. XeLaTeX 은 TikZ 그래픽을 PDF special 로 내보내는데
%   dvisvgm 이 그것을 받으려면 Ghostscript < 10.01 이 필요하다. 이 PC 에는
%   10.07 이 깔려 있어 그 경로가 막힌다. latex -> DVI -> dvisvgm 이 유일하게
%   동작하는 조합이다.

\usepackage[hangul]{kotex}
\usepackage{tikz}
\usepackage{amsmath}
\usepackage{xcolor}
\usetikzlibrary{arrows.meta, positioning, calc, fit, backgrounds, decorations.pathmorphing}

% ---- 색 --------------------------------------------------------------------
% 강조는 하나만. 나머지는 검정. 흑백 인쇄에서도 읽혀야 한다.
\definecolor{ink}{HTML}{1A1A1A}
\definecolor{soft}{HTML}{666666}
\definecolor{accent}{HTML}{7C3AED}
\definecolor{warn}{HTML}{B91C1C}

% 계층 색 — Simulink 모델의 블록 색과 같은 규칙을 쓴다
\definecolor{guid}{HTML}{1D4ED8}    % 유도
\definecolor{ctrl}{HTML}{EA580C}    % 제어
\definecolor{thr}{HTML}{D97706}     % 추진기
\definecolor{plant}{HTML}{16A34A}   % 운동모델
\definecolor{nav}{HTML}{7C3AED}     % 항법 · 신호처리
\definecolor{pathc}{HTML}{0D9488}   % 계획 경로 · 항적
\definecolor{errc}{HTML}{C2410C}    % 오차
\definecolor{flow}{HTML}{0369A1}    % 조류 · 바람

% ---- 선과 화살촉 -----------------------------------------------------------
\tikzset{
  sig/.style   = {ink, line width=0.5pt, -{Stealth[length=4pt, width=3pt]}},
  wire/.style  = {ink, line width=0.5pt},
  fb/.style    = {ink, line width=0.5pt, densely dashed,
                  -{Stealth[length=4pt, width=3pt]}},
  asig/.style  = {accent, line width=0.5pt, densely dashed,
                  -{Stealth[length=4pt, width=3pt]}},
  % 블록 — 흰 바탕, 직각, 검은 테두리. 색은 테두리에만 준다
  blk/.style   = {draw=ink, line width=0.55pt, fill=white,
                  minimum width=18mm, minimum height=8.5mm, inner sep=2pt, font=\small},
  gblk/.style  = {blk, draw=guid},
  cblk/.style  = {blk, draw=ctrl},
  tblk/.style  = {blk, draw=thr},
  pblk/.style  = {blk, draw=plant},
  nblk/.style  = {blk, draw=nav},
  % 합산점 · 분기점
  sum/.style   = {draw=ink, line width=0.55pt, fill=white, circle,
                  inner sep=0pt, minimum size=4.6mm, font=\scriptsize},
  dot/.style   = {ink, fill=ink, circle, inner sep=0pt, minimum size=1.5mm},
  % 글자
  sname/.style = {font=\small, inner sep=1.5pt},
  note/.style  = {font=\scriptsize, soft, inner sep=1.5pt},
  wnote/.style = {font=\scriptsize, warn, inner sep=1.5pt},
  anote/.style = {font=\scriptsize, accent, inner sep=1.5pt},
  axis/.style  = {ink, line width=0.5pt},
  tick/.style  = {font=\scriptsize, soft},
  % 묶음 상자 — 서브시스템 표시
  grp/.style   = {draw=soft, line width=0.4pt, densely dotted, rounded corners=1mm,
                  inner sep=3mm},
}
   (assets/src/ 안에서)
% 빌드:     bash _tools/tikz2svg.sh assets/src/w02-node-topic.tex assets/w02-node-topic.svg
%
% 왜 TikZ 인가
%   손으로 SVG 좌표를 쓰면 화살촉·선 굵기·글자 크기가 그림마다 미묘하게 달라진다.
%   그것이 "자료가 어설퍼 보인다"의 정체다. 스타일을 한 번만 정하고 상대좌표로
%   배치하면 좌표를 고쳐도 그림이 무너지지 않는다.
%   수식은 본문과 같은 TeX 글꼴로 나온다.
%
% 한글
%   ko.TeX(DVI 경로)를 쓴다. XeLaTeX 은 TikZ 그래픽을 PDF special 로 내보내는데
%   dvisvgm 이 그것을 받으려면 Ghostscript < 10.01 이 필요하다. 이 PC 에는
%   10.07 이 깔려 있어 그 경로가 막힌다. latex -> DVI -> dvisvgm 이 유일하게
%   동작하는 조합이다.

\usepackage[hangul]{kotex}
\usepackage{tikz}
\usepackage{amsmath}
\usepackage{xcolor}
\usetikzlibrary{arrows.meta, positioning, calc, fit, backgrounds, decorations.pathmorphing}

% ---- 색 --------------------------------------------------------------------
% 강조는 하나만. 나머지는 검정. 흑백 인쇄에서도 읽혀야 한다.
\definecolor{ink}{HTML}{1A1A1A}
\definecolor{soft}{HTML}{666666}
\definecolor{accent}{HTML}{7C3AED}
\definecolor{warn}{HTML}{B91C1C}

% 계층 색 — Simulink 모델의 블록 색과 같은 규칙을 쓴다
\definecolor{guid}{HTML}{1D4ED8}    % 유도
\definecolor{ctrl}{HTML}{EA580C}    % 제어
\definecolor{thr}{HTML}{D97706}     % 추진기
\definecolor{plant}{HTML}{16A34A}   % 운동모델
\definecolor{nav}{HTML}{7C3AED}     % 항법 · 신호처리
\definecolor{pathc}{HTML}{0D9488}   % 계획 경로 · 항적
\definecolor{errc}{HTML}{C2410C}    % 오차
\definecolor{flow}{HTML}{0369A1}    % 조류 · 바람

% ---- 선과 화살촉 -----------------------------------------------------------
\tikzset{
  sig/.style   = {ink, line width=0.5pt, -{Stealth[length=4pt, width=3pt]}},
  wire/.style  = {ink, line width=0.5pt},
  fb/.style    = {ink, line width=0.5pt, densely dashed,
                  -{Stealth[length=4pt, width=3pt]}},
  asig/.style  = {accent, line width=0.5pt, densely dashed,
                  -{Stealth[length=4pt, width=3pt]}},
  % 블록 — 흰 바탕, 직각, 검은 테두리. 색은 테두리에만 준다
  blk/.style   = {draw=ink, line width=0.55pt, fill=white,
                  minimum width=18mm, minimum height=8.5mm, inner sep=2pt, font=\small},
  gblk/.style  = {blk, draw=guid},
  cblk/.style  = {blk, draw=ctrl},
  tblk/.style  = {blk, draw=thr},
  pblk/.style  = {blk, draw=plant},
  nblk/.style  = {blk, draw=nav},
  % 합산점 · 분기점
  sum/.style   = {draw=ink, line width=0.55pt, fill=white, circle,
                  inner sep=0pt, minimum size=4.6mm, font=\scriptsize},
  dot/.style   = {ink, fill=ink, circle, inner sep=0pt, minimum size=1.5mm},
  % 글자
  sname/.style = {font=\small, inner sep=1.5pt},
  note/.style  = {font=\scriptsize, soft, inner sep=1.5pt},
  wnote/.style = {font=\scriptsize, warn, inner sep=1.5pt},
  anote/.style = {font=\scriptsize, accent, inner sep=1.5pt},
  axis/.style  = {ink, line width=0.5pt},
  tick/.style  = {font=\scriptsize, soft},
  % 묶음 상자 — 서브시스템 표시
  grp/.style   = {draw=soft, line width=0.4pt, densely dotted, rounded corners=1mm,
                  inner sep=3mm},
}
   (assets/src/ 안에서)
% 빌드:     bash _tools/tikz2svg.sh assets/src/w02-node-topic.tex assets/w02-node-topic.svg
%
% 왜 TikZ 인가
%   손으로 SVG 좌표를 쓰면 화살촉·선 굵기·글자 크기가 그림마다 미묘하게 달라진다.
%   그것이 "자료가 어설퍼 보인다"의 정체다. 스타일을 한 번만 정하고 상대좌표로
%   배치하면 좌표를 고쳐도 그림이 무너지지 않는다.
%   수식은 본문과 같은 TeX 글꼴로 나온다.
%
% 한글
%   ko.TeX(DVI 경로)를 쓴다. XeLaTeX 은 TikZ 그래픽을 PDF special 로 내보내는데
%   dvisvgm 이 그것을 받으려면 Ghostscript < 10.01 이 필요하다. 이 PC 에는
%   10.07 이 깔려 있어 그 경로가 막힌다. latex -> DVI -> dvisvgm 이 유일하게
%   동작하는 조합이다.

\usepackage[hangul]{kotex}
\usepackage{tikz}
\usepackage{amsmath}
\usepackage{xcolor}
\usetikzlibrary{arrows.meta, positioning, calc, fit, backgrounds, decorations.pathmorphing}

% ---- 색 --------------------------------------------------------------------
% 강조는 하나만. 나머지는 검정. 흑백 인쇄에서도 읽혀야 한다.
\definecolor{ink}{HTML}{1A1A1A}
\definecolor{soft}{HTML}{666666}
\definecolor{accent}{HTML}{7C3AED}
\definecolor{warn}{HTML}{B91C1C}

% 계층 색 — Simulink 모델의 블록 색과 같은 규칙을 쓴다
\definecolor{guid}{HTML}{1D4ED8}    % 유도
\definecolor{ctrl}{HTML}{EA580C}    % 제어
\definecolor{thr}{HTML}{D97706}     % 추진기
\definecolor{plant}{HTML}{16A34A}   % 운동모델
\definecolor{nav}{HTML}{7C3AED}     % 항법 · 신호처리
\definecolor{pathc}{HTML}{0D9488}   % 계획 경로 · 항적
\definecolor{errc}{HTML}{C2410C}    % 오차
\definecolor{flow}{HTML}{0369A1}    % 조류 · 바람

% ---- 선과 화살촉 -----------------------------------------------------------
\tikzset{
  sig/.style   = {ink, line width=0.5pt, -{Stealth[length=4pt, width=3pt]}},
  wire/.style  = {ink, line width=0.5pt},
  fb/.style    = {ink, line width=0.5pt, densely dashed,
                  -{Stealth[length=4pt, width=3pt]}},
  asig/.style  = {accent, line width=0.5pt, densely dashed,
                  -{Stealth[length=4pt, width=3pt]}},
  % 블록 — 흰 바탕, 직각, 검은 테두리. 색은 테두리에만 준다
  blk/.style   = {draw=ink, line width=0.55pt, fill=white,
                  minimum width=18mm, minimum height=8.5mm, inner sep=2pt, font=\small},
  gblk/.style  = {blk, draw=guid},
  cblk/.style  = {blk, draw=ctrl},
  tblk/.style  = {blk, draw=thr},
  pblk/.style  = {blk, draw=plant},
  nblk/.style  = {blk, draw=nav},
  % 합산점 · 분기점
  sum/.style   = {draw=ink, line width=0.55pt, fill=white, circle,
                  inner sep=0pt, minimum size=4.6mm, font=\scriptsize},
  dot/.style   = {ink, fill=ink, circle, inner sep=0pt, minimum size=1.5mm},
  % 글자
  sname/.style = {font=\small, inner sep=1.5pt},
  note/.style  = {font=\scriptsize, soft, inner sep=1.5pt},
  wnote/.style = {font=\scriptsize, warn, inner sep=1.5pt},
  anote/.style = {font=\scriptsize, accent, inner sep=1.5pt},
  axis/.style  = {ink, line width=0.5pt},
  tick/.style  = {font=\scriptsize, soft},
  % 묶음 상자 — 서브시스템 표시
  grp/.style   = {draw=soft, line width=0.4pt, densely dotted, rounded corners=1mm,
                  inner sep=3mm},
}


\begin{document}
\begin{tikzpicture}[x=1mm, y=1mm]

\pgfmathsetmacro{\PS}{3.4}                       % 1 m = 3.4 mm
\newcommand{\TN}[2]{({(#2)*\PS},{(#1)*\PS})}     % #1 = 부두 방향 t, #2 = 부두 바깥 n

% --- 배 한 척 --------------------------------------------------------------
%   #1 = t [m]   #2 = n [m]   #3 = 선수각 psi [deg]   #4 = 외곽선 스타일
\newcommand{\boat}[4]{%
  \begin{scope}[shift={({(#2)*\PS},{(#1)*\PS})}, rotate={90-(#3)}]
    \draw[#4] (-5.8,-3.6) -- (-5.8,3.6) -- (2.2,3.6) -- (7.0,0) -- (2.2,-3.6) -- cycle;
    \draw[#4, line width=0.3pt] (-5.8,0) -- (7.0,0);
  \end{scope}}

% --- 배가 가는 방향 화살표 (선체 기준) --------------------------------------
%   #1 = t   #2 = n   #3 = psi [deg]   #4 = 선체축 각 (0 = 선수, -90 = 우현)   #5 = 색
\newcommand{\boatvec}[5]{%
  \begin{scope}[shift={({(#2)*\PS},{(#1)*\PS})}, rotate={90-(#3)}]
    \draw[#5, line width=1.1pt, -{Stealth[length=4pt,width=3pt]}]
         (#4:10mm) -- (#4:22mm);
  \end{scope}}

% --- 부두 -------------------------------------------------------------------
\newcommand{\quay}{%
  \fill[black!7]  \TN{11}{-2.6} -- \TN{-11}{-2.6} -- \TN{-11}{-9} -- \TN{11}{-9} -- cycle;
  \draw[soft, line width=0.5pt]
                  \TN{11}{-2.6} -- \TN{-11}{-2.6} -- \TN{-11}{-9} -- \TN{11}{-9} -- cycle;
  \begin{scope}
    \clip \TN{11}{-2.6} -- \TN{-11}{-2.6} -- \TN{-11}{-9} -- \TN{11}{-9} -- cycle;
    \foreach \k in {-13,-11,...,11}{%
        \draw[soft, line width=0.25pt] \TN{\k}{-1.5} -- \TN{\k+1.7}{-9};}
  \end{scope}
  \node[font=\scriptsize, soft] at \TN{0}{-6.2} {부두};
  % 선석 · 접근점 마커
  \node[rectangle, draw=warn, fill=white, line width=0.7pt, inner sep=1.5pt] at \TN{0}{0} {};
  \node[circle, fill=plant, inner sep=1.3pt] at \TN{0}{15} {};}

% --- 반경 (배 위에 그려 가려지지 않게 한다) ---------------------------------
\newcommand{\radii}{%
  \draw[plant, line width=0.4pt, densely dashed] \TN{0}{15} circle ({3.0*\PS});
  \draw[warn,  line width=0.4pt, densely dashed] \TN{0}{0}  circle ({1.5*\PS});}

% --- 접근 경로 (단계 0) -----------------------------------------------------
\newcommand{\approachpath}{%
  \draw[guid, line width=1.5pt, -{Stealth[length=5pt,width=4pt]}]
        \TN{-13}{1} .. controls \TN{-8.5}{5} and \TN{-4}{11} .. \TN{0}{12.6};}

% ===========================================================================
% (가) 옆걸음 접안 — psi_f = psi_t
% ===========================================================================
\begin{scope}[shift={(0,0)}]
  \node[font=\small\bfseries, anchor=north west] at (-50,48)
       {(가) 옆걸음 접안 — $\psi_f = \psi_t$ : 부두와 나란히 서서 옆으로 붙는다};

  \quay
  \approachpath

  % 단계 0 — 경로를 따라온다. 선수각은 경로 방향
  \boat{-9.2}{3.6}{56}{soft, line width=0.5pt, fill=white, fill opacity=0.88}
  \node[font=\scriptsize, guid, anchor=center] at \TN{-11.8}{10.5}
       {\textbf{0} 접근 — 접근점까지};

  % 단계 1 — 접근점에서 선수각을 psi_t 로 돌린다
  \boat{0}{15}{20}{soft, line width=0.4pt, densely dashed}      % 돌기 전
  \boat{0}{15}{0}{ink, line width=0.6pt, fill=white}            % 돌고 난 뒤
  \draw[errc, line width=0.5pt, -{Stealth[length=3pt,width=2pt]}]
        ({15*\PS},{15}) arc (90:40:15mm);
  \node[font=\scriptsize, errc, anchor=center] at \TN{4.6}{17.8} {$\psi \to \psi_t$};
  \node[font=\scriptsize, ctrl, anchor=center] at \TN{-5.2}{19.6}
       {\textbf{1} 선수각 정렬 — $|\psi - \psi_f| < 5^{\circ}$};

  % 단계 2 — 옆걸음. 배는 북을 보고 서쪽으로 미끄러진다
  \boat{0}{7.5}{0}{ink, line width=0.6pt, fill=white, fill opacity=0.88}
  \boatvec{0}{7.5}{0}{90}{warn}
  \node[font=\scriptsize, warn, anchor=center] at \TN{5.4}{8.0}
       {\textbf{2} 접안 — 선체 $y$ 방향};

  % 단계 3 — 선석에 선다
  \boat{0}{0}{0}{ink, line width=0.7pt, fill=white, fill opacity=0.88}
  \node[font=\scriptsize, warn, anchor=center] at \TN{-5.0}{4.0}
       {\textbf{3} 정지 — $R_{stop}$ 안};
  \radii

  % 어느 추진기가 일하는가
  \node[anchor=north west, font=\scriptsize, text width=62mm] at (-50,-52)
       {\textbf{일하는 추진기} \\[2pt]
        배가 북을 보고 서쪽으로 가므로 \textbf{옆으로 미는 힘} $Y$ 가 필요하다.
        $Y$ 를 낼 수 있는 것은 \textcolor{thr}{선수추진기 $F_B$} 하나뿐이다 \\[2pt]
        그런데 $F_B$ 는 뱃머리에 있어 혼자 밀면 배가 돌아간다.
        그 돌림을 \textcolor{thr}{좌우 추진기 $F_L,\ F_R$} 이 반대로 밀어 지운다 \\[2pt]
        $\;\to\;$ 세 추진기가 모두 일한다. 선수추진기가 없으면 이 전략은 쓸 수 없다};
\end{scope}

% ===========================================================================
% (나) 전진 접안 — psi_f = psi_t - s*90도
% ===========================================================================
\begin{scope}[shift={(168,0)}]
  \node[font=\small\bfseries, anchor=north west] at (-50,48)
       {(나) 전진 접안 — $\psi_f = \psi_t - s\,90^{\circ}$ : 부두를 보고 앞으로 붙는다};

  \quay
  \approachpath

  \boat{-9.2}{3.6}{56}{soft, line width=0.5pt, fill=white, fill opacity=0.88}
  \node[font=\scriptsize, guid, anchor=center] at \TN{-11.8}{10.5}
       {\textbf{0} 접근 — 접근점까지};

  % 단계 1 — 접근점에서 부두 쪽으로 90도 돌린다
  \boat{0}{15}{20}{soft, line width=0.4pt, densely dashed}
  \boat{0}{15}{-90}{ink, line width=0.6pt, fill=white, fill opacity=0.88}
  \draw[errc, line width=0.5pt, -{Stealth[length=3pt,width=2pt]}]
        ({15*\PS},{15}) arc (90:180:15mm);
  \node[font=\scriptsize, errc, anchor=center] at \TN{4.8}{18.2} {$\psi \to \psi_t - 90^{\circ}$};
  \node[font=\scriptsize, ctrl, anchor=center] at \TN{-5.2}{19.6}
       {\textbf{1} 선수각 정렬 — $|\psi - \psi_f| < 5^{\circ}$};

  % 단계 2 — 전진
  \boat{0}{7.5}{-90}{ink, line width=0.6pt, fill=white, fill opacity=0.88}
  \boatvec{0}{7.5}{-90}{0}{warn}
  \node[font=\scriptsize, warn, anchor=center] at \TN{5.4}{8.0}
       {\textbf{2} 접안 — 선체 $x$ 방향};

  \boat{0}{0}{-90}{ink, line width=0.7pt, fill=white, fill opacity=0.88}
  \node[font=\scriptsize, warn, anchor=center] at \TN{-5.0}{4.0}
       {\textbf{3} 정지 — $R_{stop}$ 안};
  \radii

  \node[anchor=north west, font=\scriptsize, text width=62mm] at (-50,-52)
       {\textbf{일하는 추진기} \\[2pt]
        배가 가려는 쪽이 곧 뱃머리 쪽이므로 \textbf{앞으로 미는 힘} $X$ 만 있으면 된다.
        \textcolor{thr}{좌우 추진기 $F_L,\ F_R$} 이 함께 민다 \\[2pt]
        선수추진기는 거의 쉰다. 대신 마지막에 멈추려면 두 추진기를
        \textbf{뒤로} 돌려야 하고, 여기서 역 프로펠러 효율 문제가 나온다 \\[2pt]
        $\;\to\;$ 선수추진기 없는 배도 접안할 수 있다. 대신 멈추기가 어렵다};
\end{scope}

% ===========================================================================
% 아래 띠 — 두 전략이 갈리는 한 줄
% ===========================================================================
\draw[wire, soft] (-50,-98) -- (236,-98);

\node[anchor=north west, font=\scriptsize, text width=300mm] at (-50,-103)
     {\textbf{갈리는 곳은 딱 한 줄이다} — 최종 선수각 $\psi_f$ 를 무엇으로 두는가.
      나머지(접근점, 단계 전이 조건, 제어기, 배분행렬)는 둘이 똑같다 \\[4pt]
      $\psi_f = \psi_t$ (옆걸음) \qquad 또는 \qquad
      $\psi_f = \psi_t - s\,\dfrac{\pi}{2}$ (전진) \\[4pt]
      $\psi_t = 0,\ s = +1$ 을 넣으면 옆걸음은 $\psi_f = 0^{\circ}$ (북쪽을 봄),
      전진은 $\psi_f = -90^{\circ}$ (서쪽 = 부두 쪽을 봄)};

\end{tikzpicture}
\end{document}
